\documentclass[11pt, a4paper]{article}
% ── Shared preamble for all RTOS lab documents ────────────────────────────────
% Usage: % ── Shared preamble for all RTOS lab documents ────────────────────────────────
% Usage: % ── Shared preamble for all RTOS lab documents ────────────────────────────────
% Usage: \input{../preamble}  (from each labNN/ directory)
% ──────────────────────────────────────────────────────────────────────────────

% ── Packages ──────────────────────────────────────────────────────────────────
\usepackage[a4paper, top=2.5cm, bottom=2.5cm, left=2.5cm, right=2.5cm, headheight=28pt]{geometry}
\usepackage[utf8]{inputenc}
\usepackage[T1]{fontenc}
\usepackage{lmodern}
\usepackage{booktabs}
\usepackage{array}
\usepackage{tabularx}
\usepackage{longtable}
\usepackage{multirow}
\usepackage{colortbl}
\usepackage{xcolor}
\usepackage{titlesec}
\usepackage{enumitem}
\usepackage{hyperref}
\usepackage{fancyhdr}
\usepackage{parskip}
\usepackage{listings}
\usepackage[most]{tcolorbox}
\usepackage{fontawesome5}
\usepackage{microtype}
\usepackage{graphicx}
\usepackage{amsmath}

% ── Colors (KMUTNB brand) ──────────────────────────────────────────────────────
\definecolor{kmutnbBlue}{RGB}{0,56,116}
\definecolor{kmutnbGold}{RGB}{186,146,36}
\definecolor{lightGray}{RGB}{245,245,245}
\definecolor{midGray}{RGB}{200,200,200}
\definecolor{codeGray}{RGB}{248,248,248}
\definecolor{codeFrame}{RGB}{220,220,220}
\definecolor{tipteal}{RGB}{0,105,92}
\definecolor{warnAmber}{RGB}{121,85,0}
\definecolor{checkGreen}{RGB}{27,94,32}

% ── Section formatting ─────────────────────────────────────────────────────────
\titleformat{\section}
  {\large\bfseries\color{kmutnbBlue}}
  {\thesection.}{0.5em}{}
  [\color{kmutnbGold}\titlerule]

\titleformat{\subsection}
  {\normalsize\bfseries\color{kmutnbBlue}}
  {\thesubsection.}{0.5em}{}

\titleformat{\subsubsection}
  {\small\bfseries\color{kmutnbBlue}}
  {\thesubsubsection.}{0.5em}{}

\titlespacing*{\section}{0pt}{1.4ex plus .2ex}{0.6ex plus .1ex}
\titlespacing*{\subsection}{0pt}{1.0ex plus .2ex}{0.4ex plus .1ex}

% ── Listings (C and bash) ──────────────────────────────────────────────────────
\lstdefinestyle{cstyle}{
  language        = C,
  basicstyle      = \ttfamily\small,
  keywordstyle    = \bfseries\color{kmutnbBlue},
  commentstyle    = \itshape\color{gray},
  stringstyle     = \color{tipteal},
  numberstyle     = \tiny\color{midGray},
  numbers         = left,
  numbersep       = 6pt,
  stepnumber      = 1,
  backgroundcolor = \color{codeGray},
  frame           = single,
  framerule       = 0.4pt,
  rulecolor       = \color{codeFrame},
  breaklines      = true,
  showstringspaces= false,
  tabsize         = 4,
  xleftmargin     = 1.5em,
  xrightmargin    = 0.5em,
  aboveskip       = 0.8em,
  belowskip       = 0.8em,
}

\lstdefinestyle{bashstyle}{
  language        = bash,
  basicstyle      = \ttfamily\small,
  keywordstyle    = \color{kmutnbBlue},
  commentstyle    = \itshape\color{gray},
  backgroundcolor = \color{black!5},
  frame           = single,
  framerule       = 0.4pt,
  rulecolor       = \color{codeFrame},
  breaklines      = true,
  showstringspaces= false,
  xleftmargin     = 1.5em,
  xrightmargin    = 0.5em,
  aboveskip       = 0.6em,
  belowskip       = 0.6em,
}

\lstdefinestyle{textstyle}{
  basicstyle      = \ttfamily\footnotesize,
  backgroundcolor = \color{black!4},
  frame           = single,
  framerule       = 0.4pt,
  rulecolor       = \color{codeFrame},
  breaklines      = true,
  xleftmargin     = 1.5em,
  xrightmargin    = 0.5em,
  aboveskip       = 0.6em,
  belowskip       = 0.6em,
}

% ── Callout boxes ──────────────────────────────────────────────────────────────
\tcbuselibrary{skins, breakable}

\newtcolorbox{tipbox}[1][]{
  enhanced, breakable,
  colback=tipteal!8, colframe=tipteal, coltitle=white,
  fonttitle=\bfseries\small, title={\faLightbulb\quad Tip#1},
  left=6pt, right=6pt, top=4pt, bottom=4pt,
  boxrule=0.8pt, arc=3pt,
}

\newtcolorbox{warnbox}[1][]{
  enhanced, breakable,
  colback=warnAmber!8, colframe=kmutnbGold, coltitle=white,
  fonttitle=\bfseries\small, title={\faExclamationTriangle\quad Note#1},
  left=6pt, right=6pt, top=4pt, bottom=4pt,
  boxrule=0.8pt, arc=3pt,
}

\newtcolorbox{checkpointbox}[1][]{
  enhanced,
  colback=kmutnbBlue!6, colframe=kmutnbBlue, coltitle=white,
  fonttitle=\bfseries\small, title={\faCheckCircle\quad Checkpoint#1},
  left=6pt, right=6pt, top=4pt, bottom=4pt,
  boxrule=0.8pt, arc=3pt,
}

\newtcolorbox{questionbox}[1][]{
  enhanced, breakable,
  colback=kmutnbGold!10, colframe=kmutnbGold, coltitle=white,
  fonttitle=\bfseries\small, title={\faQuestion\quad Question#1},
  left=6pt, right=6pt, top=4pt, bottom=4pt,
  boxrule=0.8pt, arc=3pt,
}

% ── Checkbox list ──────────────────────────────────────────────────────────────
\newlist{checklist}{itemize}{1}
\setlist[checklist]{label=\faSquare[regular], leftmargin=1.5em, noitemsep}

% ── Misc helpers ───────────────────────────────────────────────────────────────
\newcommand{\file}[1]{\texttt{\small #1}}
\newcommand{\api}[1]{\texttt{\small #1}}

% ── Title block macro ─────────────────────────────────────────────────────────
% Usage: \labtitle{03}{Aperiodic Servers \& Producer--Consumer}{2}{CLO 1 \& CLO 2}
\newcommand{\labtitle}[4]{%
  \begin{center}
    {\color{kmutnbBlue}\rule{\linewidth}{2pt}}\\[6pt]
    {\LARGE\bfseries\color{kmutnbBlue} Laboratory #1}\\[4pt]
    {\large\bfseries #2}\\[2pt]
    {\normalsize Real-Time Operating Systems \enspace\textbullet\enspace
     M.Eng.\ ECE \enspace\textbullet\enspace KMUTNB}\\[8pt]
    {\color{kmutnbGold}\rule{\linewidth}{1pt}}
  \end{center}
  \vspace{0.3cm}
  \begin{tabular}{@{}p{3.8cm} p{10cm}@{}}
    \textbf{Platform}       & NXP FRDM-MCXN236 (ARM Cortex-M33 @ 150\,MHz) \\[2pt]
    \textbf{RTOS}           & FreeRTOS v10.6 \\[2pt]
    \textbf{Toolchain}      & ARM GNU Toolchain (\texttt{arm-none-eabi-gcc}) + Make \\[2pt]
    \textbf{Estimated time} & #3 hours \\[2pt]
    \textbf{CLO mapping}    & #4 \\
  \end{tabular}
  \vspace{0.4cm}
}
  (from each labNN/ directory)
% ──────────────────────────────────────────────────────────────────────────────

% ── Packages ──────────────────────────────────────────────────────────────────
\usepackage[a4paper, top=2.5cm, bottom=2.5cm, left=2.5cm, right=2.5cm, headheight=28pt]{geometry}
\usepackage[utf8]{inputenc}
\usepackage[T1]{fontenc}
\usepackage{lmodern}
\usepackage{booktabs}
\usepackage{array}
\usepackage{tabularx}
\usepackage{longtable}
\usepackage{multirow}
\usepackage{colortbl}
\usepackage{xcolor}
\usepackage{titlesec}
\usepackage{enumitem}
\usepackage{hyperref}
\usepackage{fancyhdr}
\usepackage{parskip}
\usepackage{listings}
\usepackage[most]{tcolorbox}
\usepackage{fontawesome5}
\usepackage{microtype}
\usepackage{graphicx}
\usepackage{amsmath}

% ── Colors (KMUTNB brand) ──────────────────────────────────────────────────────
\definecolor{kmutnbBlue}{RGB}{0,56,116}
\definecolor{kmutnbGold}{RGB}{186,146,36}
\definecolor{lightGray}{RGB}{245,245,245}
\definecolor{midGray}{RGB}{200,200,200}
\definecolor{codeGray}{RGB}{248,248,248}
\definecolor{codeFrame}{RGB}{220,220,220}
\definecolor{tipteal}{RGB}{0,105,92}
\definecolor{warnAmber}{RGB}{121,85,0}
\definecolor{checkGreen}{RGB}{27,94,32}

% ── Section formatting ─────────────────────────────────────────────────────────
\titleformat{\section}
  {\large\bfseries\color{kmutnbBlue}}
  {\thesection.}{0.5em}{}
  [\color{kmutnbGold}\titlerule]

\titleformat{\subsection}
  {\normalsize\bfseries\color{kmutnbBlue}}
  {\thesubsection.}{0.5em}{}

\titleformat{\subsubsection}
  {\small\bfseries\color{kmutnbBlue}}
  {\thesubsubsection.}{0.5em}{}

\titlespacing*{\section}{0pt}{1.4ex plus .2ex}{0.6ex plus .1ex}
\titlespacing*{\subsection}{0pt}{1.0ex plus .2ex}{0.4ex plus .1ex}

% ── Listings (C and bash) ──────────────────────────────────────────────────────
\lstdefinestyle{cstyle}{
  language        = C,
  basicstyle      = \ttfamily\small,
  keywordstyle    = \bfseries\color{kmutnbBlue},
  commentstyle    = \itshape\color{gray},
  stringstyle     = \color{tipteal},
  numberstyle     = \tiny\color{midGray},
  numbers         = left,
  numbersep       = 6pt,
  stepnumber      = 1,
  backgroundcolor = \color{codeGray},
  frame           = single,
  framerule       = 0.4pt,
  rulecolor       = \color{codeFrame},
  breaklines      = true,
  showstringspaces= false,
  tabsize         = 4,
  xleftmargin     = 1.5em,
  xrightmargin    = 0.5em,
  aboveskip       = 0.8em,
  belowskip       = 0.8em,
}

\lstdefinestyle{bashstyle}{
  language        = bash,
  basicstyle      = \ttfamily\small,
  keywordstyle    = \color{kmutnbBlue},
  commentstyle    = \itshape\color{gray},
  backgroundcolor = \color{black!5},
  frame           = single,
  framerule       = 0.4pt,
  rulecolor       = \color{codeFrame},
  breaklines      = true,
  showstringspaces= false,
  xleftmargin     = 1.5em,
  xrightmargin    = 0.5em,
  aboveskip       = 0.6em,
  belowskip       = 0.6em,
}

\lstdefinestyle{textstyle}{
  basicstyle      = \ttfamily\footnotesize,
  backgroundcolor = \color{black!4},
  frame           = single,
  framerule       = 0.4pt,
  rulecolor       = \color{codeFrame},
  breaklines      = true,
  xleftmargin     = 1.5em,
  xrightmargin    = 0.5em,
  aboveskip       = 0.6em,
  belowskip       = 0.6em,
}

% ── Callout boxes ──────────────────────────────────────────────────────────────
\tcbuselibrary{skins, breakable}

\newtcolorbox{tipbox}[1][]{
  enhanced, breakable,
  colback=tipteal!8, colframe=tipteal, coltitle=white,
  fonttitle=\bfseries\small, title={\faLightbulb\quad Tip#1},
  left=6pt, right=6pt, top=4pt, bottom=4pt,
  boxrule=0.8pt, arc=3pt,
}

\newtcolorbox{warnbox}[1][]{
  enhanced, breakable,
  colback=warnAmber!8, colframe=kmutnbGold, coltitle=white,
  fonttitle=\bfseries\small, title={\faExclamationTriangle\quad Note#1},
  left=6pt, right=6pt, top=4pt, bottom=4pt,
  boxrule=0.8pt, arc=3pt,
}

\newtcolorbox{checkpointbox}[1][]{
  enhanced,
  colback=kmutnbBlue!6, colframe=kmutnbBlue, coltitle=white,
  fonttitle=\bfseries\small, title={\faCheckCircle\quad Checkpoint#1},
  left=6pt, right=6pt, top=4pt, bottom=4pt,
  boxrule=0.8pt, arc=3pt,
}

\newtcolorbox{questionbox}[1][]{
  enhanced, breakable,
  colback=kmutnbGold!10, colframe=kmutnbGold, coltitle=white,
  fonttitle=\bfseries\small, title={\faQuestion\quad Question#1},
  left=6pt, right=6pt, top=4pt, bottom=4pt,
  boxrule=0.8pt, arc=3pt,
}

% ── Checkbox list ──────────────────────────────────────────────────────────────
\newlist{checklist}{itemize}{1}
\setlist[checklist]{label=\faSquare[regular], leftmargin=1.5em, noitemsep}

% ── Misc helpers ───────────────────────────────────────────────────────────────
\newcommand{\file}[1]{\texttt{\small #1}}
\newcommand{\api}[1]{\texttt{\small #1}}

% ── Title block macro ─────────────────────────────────────────────────────────
% Usage: \labtitle{03}{Aperiodic Servers \& Producer--Consumer}{2}{CLO 1 \& CLO 2}
\newcommand{\labtitle}[4]{%
  \begin{center}
    {\color{kmutnbBlue}\rule{\linewidth}{2pt}}\\[6pt]
    {\LARGE\bfseries\color{kmutnbBlue} Laboratory #1}\\[4pt]
    {\large\bfseries #2}\\[2pt]
    {\normalsize Real-Time Operating Systems \enspace\textbullet\enspace
     M.Eng.\ ECE \enspace\textbullet\enspace KMUTNB}\\[8pt]
    {\color{kmutnbGold}\rule{\linewidth}{1pt}}
  \end{center}
  \vspace{0.3cm}
  \begin{tabular}{@{}p{3.8cm} p{10cm}@{}}
    \textbf{Platform}       & NXP FRDM-MCXN236 (ARM Cortex-M33 @ 150\,MHz) \\[2pt]
    \textbf{RTOS}           & FreeRTOS v10.6 \\[2pt]
    \textbf{Toolchain}      & ARM GNU Toolchain (\texttt{arm-none-eabi-gcc}) + Make \\[2pt]
    \textbf{Estimated time} & #3 hours \\[2pt]
    \textbf{CLO mapping}    & #4 \\
  \end{tabular}
  \vspace{0.4cm}
}
  (from each labNN/ directory)
% ──────────────────────────────────────────────────────────────────────────────

% ── Packages ──────────────────────────────────────────────────────────────────
\usepackage[a4paper, top=2.5cm, bottom=2.5cm, left=2.5cm, right=2.5cm, headheight=28pt]{geometry}
\usepackage[utf8]{inputenc}
\usepackage[T1]{fontenc}
\usepackage{lmodern}
\usepackage{booktabs}
\usepackage{array}
\usepackage{tabularx}
\usepackage{longtable}
\usepackage{multirow}
\usepackage{colortbl}
\usepackage{xcolor}
\usepackage{titlesec}
\usepackage{enumitem}
\usepackage{hyperref}
\usepackage{fancyhdr}
\usepackage{parskip}
\usepackage{listings}
\usepackage[most]{tcolorbox}
\usepackage{fontawesome5}
\usepackage{microtype}
\usepackage{graphicx}
\usepackage{amsmath}

% ── Colors (KMUTNB brand) ──────────────────────────────────────────────────────
\definecolor{kmutnbBlue}{RGB}{0,56,116}
\definecolor{kmutnbGold}{RGB}{186,146,36}
\definecolor{lightGray}{RGB}{245,245,245}
\definecolor{midGray}{RGB}{200,200,200}
\definecolor{codeGray}{RGB}{248,248,248}
\definecolor{codeFrame}{RGB}{220,220,220}
\definecolor{tipteal}{RGB}{0,105,92}
\definecolor{warnAmber}{RGB}{121,85,0}
\definecolor{checkGreen}{RGB}{27,94,32}

% ── Section formatting ─────────────────────────────────────────────────────────
\titleformat{\section}
  {\large\bfseries\color{kmutnbBlue}}
  {\thesection.}{0.5em}{}
  [\color{kmutnbGold}\titlerule]

\titleformat{\subsection}
  {\normalsize\bfseries\color{kmutnbBlue}}
  {\thesubsection.}{0.5em}{}

\titleformat{\subsubsection}
  {\small\bfseries\color{kmutnbBlue}}
  {\thesubsubsection.}{0.5em}{}

\titlespacing*{\section}{0pt}{1.4ex plus .2ex}{0.6ex plus .1ex}
\titlespacing*{\subsection}{0pt}{1.0ex plus .2ex}{0.4ex plus .1ex}

% ── Listings (C and bash) ──────────────────────────────────────────────────────
\lstdefinestyle{cstyle}{
  language        = C,
  basicstyle      = \ttfamily\small,
  keywordstyle    = \bfseries\color{kmutnbBlue},
  commentstyle    = \itshape\color{gray},
  stringstyle     = \color{tipteal},
  numberstyle     = \tiny\color{midGray},
  numbers         = left,
  numbersep       = 6pt,
  stepnumber      = 1,
  backgroundcolor = \color{codeGray},
  frame           = single,
  framerule       = 0.4pt,
  rulecolor       = \color{codeFrame},
  breaklines      = true,
  showstringspaces= false,
  tabsize         = 4,
  xleftmargin     = 1.5em,
  xrightmargin    = 0.5em,
  aboveskip       = 0.8em,
  belowskip       = 0.8em,
}

\lstdefinestyle{bashstyle}{
  language        = bash,
  basicstyle      = \ttfamily\small,
  keywordstyle    = \color{kmutnbBlue},
  commentstyle    = \itshape\color{gray},
  backgroundcolor = \color{black!5},
  frame           = single,
  framerule       = 0.4pt,
  rulecolor       = \color{codeFrame},
  breaklines      = true,
  showstringspaces= false,
  xleftmargin     = 1.5em,
  xrightmargin    = 0.5em,
  aboveskip       = 0.6em,
  belowskip       = 0.6em,
}

\lstdefinestyle{textstyle}{
  basicstyle      = \ttfamily\footnotesize,
  backgroundcolor = \color{black!4},
  frame           = single,
  framerule       = 0.4pt,
  rulecolor       = \color{codeFrame},
  breaklines      = true,
  xleftmargin     = 1.5em,
  xrightmargin    = 0.5em,
  aboveskip       = 0.6em,
  belowskip       = 0.6em,
}

% ── Callout boxes ──────────────────────────────────────────────────────────────
\tcbuselibrary{skins, breakable}

\newtcolorbox{tipbox}[1][]{
  enhanced, breakable,
  colback=tipteal!8, colframe=tipteal, coltitle=white,
  fonttitle=\bfseries\small, title={\faLightbulb\quad Tip#1},
  left=6pt, right=6pt, top=4pt, bottom=4pt,
  boxrule=0.8pt, arc=3pt,
}

\newtcolorbox{warnbox}[1][]{
  enhanced, breakable,
  colback=warnAmber!8, colframe=kmutnbGold, coltitle=white,
  fonttitle=\bfseries\small, title={\faExclamationTriangle\quad Note#1},
  left=6pt, right=6pt, top=4pt, bottom=4pt,
  boxrule=0.8pt, arc=3pt,
}

\newtcolorbox{checkpointbox}[1][]{
  enhanced,
  colback=kmutnbBlue!6, colframe=kmutnbBlue, coltitle=white,
  fonttitle=\bfseries\small, title={\faCheckCircle\quad Checkpoint#1},
  left=6pt, right=6pt, top=4pt, bottom=4pt,
  boxrule=0.8pt, arc=3pt,
}

\newtcolorbox{questionbox}[1][]{
  enhanced, breakable,
  colback=kmutnbGold!10, colframe=kmutnbGold, coltitle=white,
  fonttitle=\bfseries\small, title={\faQuestion\quad Question#1},
  left=6pt, right=6pt, top=4pt, bottom=4pt,
  boxrule=0.8pt, arc=3pt,
}

% ── Checkbox list ──────────────────────────────────────────────────────────────
\newlist{checklist}{itemize}{1}
\setlist[checklist]{label=\faSquare[regular], leftmargin=1.5em, noitemsep}

% ── Misc helpers ───────────────────────────────────────────────────────────────
\newcommand{\file}[1]{\texttt{\small #1}}
\newcommand{\api}[1]{\texttt{\small #1}}

% ── Title block macro ─────────────────────────────────────────────────────────
% Usage: \labtitle{03}{Aperiodic Servers \& Producer--Consumer}{2}{CLO 1 \& CLO 2}
\newcommand{\labtitle}[4]{%
  \begin{center}
    {\color{kmutnbBlue}\rule{\linewidth}{2pt}}\\[6pt]
    {\LARGE\bfseries\color{kmutnbBlue} Laboratory #1}\\[4pt]
    {\large\bfseries #2}\\[2pt]
    {\normalsize Real-Time Operating Systems \enspace\textbullet\enspace
     M.Eng.\ ECE \enspace\textbullet\enspace KMUTNB}\\[8pt]
    {\color{kmutnbGold}\rule{\linewidth}{1pt}}
  \end{center}
  \vspace{0.3cm}
  \begin{tabular}{@{}p{3.8cm} p{10cm}@{}}
    \textbf{Platform}       & NXP FRDM-MCXN236 (ARM Cortex-M33 @ 150\,MHz) \\[2pt]
    \textbf{RTOS}           & FreeRTOS v10.6 \\[2pt]
    \textbf{Toolchain}      & ARM GNU Toolchain (\texttt{arm-none-eabi-gcc}) + Make \\[2pt]
    \textbf{Estimated time} & #3 hours \\[2pt]
    \textbf{CLO mapping}    & #4 \\
  \end{tabular}
  \vspace{0.4cm}
}


% ── Header / Footer ────────────────────────────────────────────────────────────
\pagestyle{fancy}
\fancyhf{}
\fancyhead[L]{\small\color{kmutnbBlue}\textbf{KMUTNB -- Faculty of Engineering}}
\fancyhead[R]{\small\color{kmutnbBlue}Dept.\ of Electrical and Computer Engineering}
\fancyfoot[L]{\small\color{kmutnbBlue}Real-Time Operating Systems -- M.Eng.\ ECE}
\fancyfoot[C]{\small\thepage}
\fancyfoot[R]{\small\color{kmutnbBlue}Lab 07}
\renewcommand{\headrulewidth}{0.4pt}
\renewcommand{\headrule}{\hbox to\headwidth{%
  \color{kmutnbGold}\leaders\hrule height\headrulewidth\hfill}}
\renewcommand{\footrulewidth}{0.4pt}
\renewcommand{\footrule}{\hbox to\headwidth{%
  \color{midGray}\leaders\hrule height\footrulewidth\hfill}}

\hypersetup{
  colorlinks = true,
  linkcolor  = kmutnbBlue,
  urlcolor   = kmutnbBlue,
  citecolor  = kmutnbBlue,
  pdftitle   = {Lab 07 -- Stack Overflow Detection \& Memory Analysis},
  pdfauthor  = {KMUTNB RTOS Course},
}

% ══════════════════════════════════════════════════════════════════════════════
\begin{document}

\labtitle{07}{Stack Overflow Detection \& Memory Analysis}{2--3}{CLO 4 \& CLO 5}

% ── Objectives ────────────────────────────────────────────────────────────────
\section{Objectives}

By completing this lab you will be able to:

\begin{enumerate}[noitemsep, leftmargin=*]
  \item \textbf{Induce} a stack overflow in a FreeRTOS task using a recursive function.
  \item \textbf{Detect} the overflow using \texttt{configCHECK\_FOR\_STACK\_OVERFLOW = 2}
        before corruption spreads.
  \item \textbf{Right-size} task stacks using \texttt{uxTaskGetStackHighWaterMark}.
  \item \textbf{Analyse} total heap usage with \texttt{xPortGetFreeHeapSize} and the
        linker map file.
  \item \textbf{Convert} one task to static allocation and verify heap usage drops.
\end{enumerate}

% ── Prerequisites ─────────────────────────────────────────────────────────────
\section{Prerequisites}

\begin{checklist}
  \item Lab 03 producer-consumer system (used as the baseline in Part~C).
  \item ARM GNU Toolchain and \texttt{pyocd} working.
  \item A text editor with hex search (for observing the 0xA5 pattern).
\end{checklist}

% ── Background ────────────────────────────────────────────────────────────────
\section{Background}

\subsection{Stack Overflow Consequences}

Every FreeRTOS task has its own stack allocated from the heap. If a task's stack grows
beyond its allocation, it silently overwrites the adjacent heap block --- which may
contain another task's TCB, stack, or queue buffer. The crash typically manifests as a
hard fault far from the actual cause.

FreeRTOS provides two overflow detection mechanisms:

\renewcommand{\arraystretch}{1.3}
\begin{tabular}{@{} c p{2cm} p{5cm} p{3cm} @{}}
  \toprule
  \textbf{Method} & \textbf{Config value} & \textbf{How it works} & \textbf{Catches} \\
  \midrule
  1 & \texttt{1} & Check SP $<$ stack base at each context switch &
    Overflows that leave SP below base \\
  2 & \texttt{2} & Fill entire stack with 0xA5 at creation; check last 20 bytes at each switch &
    Most overflows, including those that recover \\
  \bottomrule
\end{tabular}
\renewcommand{\arraystretch}{1}

% ── Project Structure ─────────────────────────────────────────────────────────
\section{Project Structure}

\begin{lstlisting}[style=textstyle]
lab07_stack_overflow/
  Makefile
  linker_MCXN236.ld
  src/
    main.c                 <- task creation
    overflow_demo.c/.h     <- Part A: intentional overflow
    stack_sizing.c/.h      <- Part B: sizing the Lab 03 pipeline
    static_alloc.c/.h      <- Part C: static allocation
    FreeRTOSConfig.h
    uart.h / uart.c
  FreeRTOS-Kernel/
\end{lstlisting}

% ── Part A ─────────────────────────────────────────────────────────────────────
\section{Part A --- Induce and Detect a Stack Overflow}

\subsection{Step 1 --- Create a Task with a Small Stack}

\begin{lstlisting}[style=cstyle]
/* overflow_demo.c */

/* Each call adds ~64 bytes to the stack */
static uint32_t recursive_sum(uint32_t n)
{
    uint32_t local_buffer[8];
    for (int i = 0; i < 8; i++) local_buffer[i] = i + n;
    if (n == 0) return 0;
    return local_buffer[0] + recursive_sum(n - 1);
}

void vOverflowTask(void *pv)
{
    for (;;) {
        uart_printf("[OVF] starting recursive_sum(60)...\r\n");
        uint32_t result = recursive_sum(60);
        uart_printf("[OVF] result = %lu (you should never see this)\r\n", result);
        vTaskDelay(pdMS_TO_TICKS(2000));
    }
}
\end{lstlisting}

In \file{main.c}, create the task with a very small stack (64 words = 256 bytes):

\begin{lstlisting}[style=cstyle]
xTaskCreate(vOverflowTask, "Overflow", 64, NULL, 2, NULL);
\end{lstlisting}

\subsection{Step 2 --- First Run: No Detection}

Build and flash with detection \textbf{disabled}:

\begin{lstlisting}[style=cstyle]
/* FreeRTOSConfig.h */
#define configCHECK_FOR_STACK_OVERFLOW  0
\end{lstlisting}

Observe the terminal. The overflow will likely cause a HardFault or garbled output, but
no useful diagnostic message.

\begin{checkpointbox}[~A1]
  UART capture showing the crash --- note that the crash is non-deterministic and may
  look different each run.
\end{checkpointbox}

\subsection{Step 3 --- Enable Overflow Detection}

\begin{lstlisting}[style=cstyle]
/* FreeRTOSConfig.h */
#define configCHECK_FOR_STACK_OVERFLOW  2
\end{lstlisting}

Implement the hook in \file{overflow\_demo.c}:

\begin{lstlisting}[style=cstyle]
void vApplicationStackOverflowHook(TaskHandle_t xTask, char *pcTaskName)
{
    taskDISABLE_INTERRUPTS();
    uart_printf("\r\n!!! STACK OVERFLOW in task: %s !!!\r\n", pcTaskName);
    for (;;) {}   /* halt so we can read the message */
}
\end{lstlisting}

Rebuild and flash.

\begin{checkpointbox}[~A2]
  UART capture showing \texttt{!!! STACK OVERFLOW in task: Overflow !!!}
\end{checkpointbox}

\begin{questionbox}[~A1]
  With \texttt{configCHECK\_FOR\_STACK\_OVERFLOW = 2}, the check only fires at context
  switch time --- not at the exact moment of overflow. Is there any scenario where the
  overflow happens, the stack ``heals'' before the switch, and the hook does NOT fire?
  When would you use Method~1 vs.\ Method~2?
\end{questionbox}

\begin{questionbox}[~A2]
  The overflow hook receives \texttt{pcTaskName}. Can you always trust this name to
  identify the crashed task? (Hint: what was overwritten just before the hook was called?)
\end{questionbox}

% ── Part B ─────────────────────────────────────────────────────────────────────
\section{Part B --- Right-Sizing Stack Depths}

\subsection{Step 1 --- Find the High-Water Mark}

Increase the overflow task's stack to 512 words so it can run safely:

\begin{lstlisting}[style=cstyle]
xTaskCreate(vOverflowTask, "Overflow", 512, NULL, 2, NULL);
\end{lstlisting}

After at least 30 seconds, read the high-water mark:

\begin{lstlisting}[style=cstyle]
void vStackReportTask(void *pv)
{
    vTaskDelay(pdMS_TO_TICKS(30000));
    uart_printf("\r\n=== Stack High-Water Mark Report ===\r\n");
    TaskHandle_t handles[] = { xOverflowHandle, xSensor1Handle, xLoggerHandle };
    const char  *names[]   = { "Overflow", "Sensor1", "Logger" };
    for (int i = 0; i < 3; i++) {
        UBaseType_t hwm = uxTaskGetStackHighWaterMark(handles[i]);
        uart_printf("  %-16s  HWM = %u words remaining\r\n", names[i], hwm);
    }
    uart_printf("Free heap: %u bytes\r\n", xPortGetFreeHeapSize());
    vTaskDelete(NULL);
}
\end{lstlisting}

\subsection{Step 2 --- Compute the Right-Sized Stack}

For each task, apply the formula:
\[
  \text{right\_sized} = (\text{allocated} - \text{HWM}) + \lceil (\text{allocated} - \text{HWM}) \times 0.20 \rceil
\]

\subsection{Step 3 --- Apply to the Lab 03 Pipeline}

\renewcommand{\arraystretch}{1.3}
\begin{tabular}{@{} p{3.5cm} c c c c @{}}
  \toprule
  \textbf{Task} & \textbf{Alloc'd (words)} & \textbf{HWM remaining} &
    \textbf{Min used} & \textbf{Right-sized (+20\%)} \\
  \midrule
  \texttt{vSensor1Task}  & & & & \\
  \texttt{vSensor2Task}  & & & & \\
  \texttt{vLoggerTask}   & & & & \\
  \texttt{vMonitorTask}  & & & & \\
  \texttt{vOverflowTask} & & & & \\
  \bottomrule
\end{tabular}
\renewcommand{\arraystretch}{1}

\begin{checkpointbox}[~B]
  Filled-in high-water-mark table from UART output.
\end{checkpointbox}

\begin{questionbox}[~B1]
  \texttt{uxTaskGetStackHighWaterMark} returns 4 for one of your tasks. What does this
  mean? Is this task safe to leave running? What is your next step?
\end{questionbox}

\begin{questionbox}[~B2]
  The formula adds 20\% headroom. What real-world factors could cause a task to use more
  stack in production than was observed in the lab? Name at least three.
\end{questionbox}

% ── Part C ─────────────────────────────────────────────────────────────────────
\section{Part C --- Heap Analysis}

\subsection{Step 1 --- Total Heap Usage}

\begin{lstlisting}[style=cstyle]
/* In main.c, before vTaskStartScheduler: */
uart_printf("Heap free before scheduler: %u bytes\r\n",
            xPortGetFreeHeapSize());
uart_printf("Heap minimum ever free:     %u bytes\r\n",
            xPortGetMinimumEverFreeHeapSize());
\end{lstlisting}

\subsection{Step 2 --- Linker Map Analysis}

Build with map output enabled:

\begin{lstlisting}[style=bashstyle]
make MAP=1
\end{lstlisting}

Open \file{lab07.map} and locate the memory layout. Fill in this table:

\renewcommand{\arraystretch}{1.3}
\begin{tabular}{@{} p{2.5cm} p{2cm} p{1.8cm} p{5.5cm} @{}}
  \toprule
  \textbf{Section} & \textbf{Size (bytes)} & \textbf{Location} & \textbf{Contains} \\
  \midrule
  \texttt{.text}       & & Flash & Code, string literals, ISR vectors \\
  \texttt{.rodata}     & & Flash & \texttt{const} data \\
  \texttt{.data}       & & SRAM  & Initialised globals \\
  \texttt{.bss}        & & SRAM  & Zero-init globals (incl.\ \texttt{ucHeap}) \\
  Stack (MSP)          & & SRAM  & Main stack (before scheduler) \\
  \bottomrule
\end{tabular}
\renewcommand{\arraystretch}{1}

\begin{checkpointbox}[~C]
  Filled-in linker map section size table.
\end{checkpointbox}

% ── Part D ─────────────────────────────────────────────────────────────────────
\section{Part D --- Static Allocation}

\subsection{Step 1 --- Convert One Task to Static Allocation}

Convert \texttt{vLoggerTask} from dynamic to static allocation:

\begin{lstlisting}[style=cstyle]
/* static_alloc.c */
static StaticTask_t xLoggerTaskBuffer;
static StackType_t  xLoggerStack[256];   /* right-sized from Part B */

TaskHandle_t xLoggerHandle = xTaskCreateStatic(
    vLoggerTask, "Logger", 256, NULL, 1,
    xLoggerStack, &xLoggerTaskBuffer
);
\end{lstlisting}

Ensure \file{FreeRTOSConfig.h} has:

\begin{lstlisting}[style=cstyle]
#define configSUPPORT_STATIC_ALLOCATION   1
#define configSUPPORT_DYNAMIC_ALLOCATION  1

void vApplicationGetIdleTaskMemory(StaticTask_t **ppxIdleTaskTCBBuffer,
                                    StackType_t  **ppxIdleTaskStackBuffer,
                                    uint32_t      *pulIdleTaskStackSize)
{
    static StaticTask_t xIdleTaskTCB;
    static StackType_t  xIdleTaskStack[configMINIMAL_STACK_SIZE];
    *ppxIdleTaskTCBBuffer   = &xIdleTaskTCB;
    *ppxIdleTaskStackBuffer = xIdleTaskStack;
    *pulIdleTaskStackSize   = configMINIMAL_STACK_SIZE;
}
\end{lstlisting}

\subsection{Step 2 --- Measure Heap Change}

Run \texttt{xPortGetFreeHeapSize()} before and after switching to static allocation.
Expected: free heap increases by approximately
\texttt{sizeof(StaticTask\_t) + 256*4} bytes.

\begin{checkpointbox}[~D]
  UART output showing free heap before and after the conversion.
\end{checkpointbox}

\begin{questionbox}[~D1]
  With \texttt{configSUPPORT\_STATIC\_ALLOCATION = 1} and
  \texttt{configSUPPORT\_DYNAMIC\_ALLOCATION = 0} (all static), what are the benefits
  for a safety-critical system? What does the linker map tell you that
  \texttt{xPortGetFreeHeapSize()} cannot?
\end{questionbox}

\begin{questionbox}[~D2]
  \texttt{vApplicationGetIdleTaskMemory} must be implemented when static allocation is
  enabled. Why does the idle task require static allocation even if all your application
  tasks use dynamic allocation?
\end{questionbox}

% ── Deliverables ───────────────────────────────────────────────────────────────
\section{Deliverables}

\renewcommand{\arraystretch}{1.3}
\begin{tabular}{@{} c p{4.5cm} p{7cm} @{}}
  \toprule
  \textbf{\#} & \textbf{Item} & \textbf{Details} \\
  \midrule
  1 & Part A1 UART capture  & Crash without detection \\
  2 & Part A2 UART capture  & Stack overflow hook message \\
  3 & Part B HWM table      & All Lab~03 tasks + overflow task \\
  4 & Part C map table      & Linker section sizes \\
  5 & Part D UART capture   & Heap usage before and after static conversion \\
  6 & Written answers       & Questions A1--A2, B1--B2, D1--D2 \\
  \bottomrule
\end{tabular}
\renewcommand{\arraystretch}{1}

% ── Key Concepts ───────────────────────────────────────────────────────────────
\section{Key Concepts Demonstrated}

\renewcommand{\arraystretch}{1.3}
\begin{tabular}{@{} p{7cm} p{6cm} @{}}
  \toprule
  \textbf{Concept} & \textbf{Where you saw it} \\
  \midrule
  Stack overflow $\to$ silent corruption       & Part A (no detection) \\
  \texttt{vApplicationStackOverflowHook}       & Part A (with detection) \\
  High-water-mark measurement and right-sizing & Part B \\
  Heap usage accounting                        & Part C \\
  Linker map section analysis                  & Part C \\
  Static task allocation (\texttt{xTaskCreateStatic}) & Part D \\
  \bottomrule
\end{tabular}
\renewcommand{\arraystretch}{1}

% ── Reference ─────────────────────────────────────────────────────────────────
\section{Reference}

\renewcommand{\arraystretch}{1.3}
\begin{tabular}{@{} p{5.5cm} p{7.5cm} @{}}
  \toprule
  \textbf{Resource} & \textbf{Relevant sections} \\
  \midrule
  Barry, \textit{Mastering the FreeRTOS Kernel} & Ch.~3 (heap schemes) \\
  FreeRTOS docs &
    \texttt{configCHECK\_FOR\_STACK\_OVERFLOW},
    \texttt{uxTaskGetStackHighWaterMark},
    \texttt{xTaskCreateStatic} \\
  NXP MCXN236 Reference Manual & \S SRAM memory map \\
  \bottomrule
\end{tabular}
\renewcommand{\arraystretch}{1}

\end{document}
